% !TEX root = ../main.tex
\documentclass[../main.tex]{subfiles}
\graphicspath{{\subfix{../images/}}}

\begin{document}

В ходе выполнения курсовой работы разработана игра в жанре bullet hell на языке C++ с использованием библиотеки Raylib. Реализованы следующие компоненты:

\begin{enumerate}
	\item Проведён анализ предметной области и выбран технологический стек (C++20, Raylib, Premake).
	\item Спроектирована объектная модель: структуры \texttt{Player}, \texttt{Projectile}, \texttt{AttackEvent}, перечисление состояний игры \texttt{GameState}. Архитектура построена на основе событийной системы, что позволяет гибко описывать уровни без изменения основного игрового цикла.
	\item Реализованы базовые механики: движение игрока с нормализацией скорости, обработка коллизий (круг–прямоугольник), система здоровья и неуязвимости (временное окно 0,5 с после попадания).
	\item Разработана система событий, включающая планировщик на основе абсолютного времени уровня и библиотеку вспомогательных функций для генерации типовых паттернов атак (\texttt{add\_beam\_events}, \texttt{spawn\_circular\_burst}, \texttt{add\_fan\_attack} и др.).
	\item Созданы два демонстрационных уровня — обучающий «Beginning» и более сложный «Way» — с различными паттернами: линейные залпы, веерные выстрелы, бомбы с взрывами, прицельные снаряды, каскадные атаки.
	\item Проведено тестирование с использованием специальных отладочных уровней (\texttt{TEST}, \texttt{TEST\_HP}, \texttt{EMPTY}).
	Оценена производительность~— частота кадров остаётся не ниже 60 FPS при 300 активных снарядах на современном оборудовании.
	\item Подготовлена инструкция по сборке и запуску проекта в среде Linux с использованием Premake5 и GCC.
\end{enumerate}

Поставленные задачи выполнены в полном объёме. Разработанная игра может служить основой для дальнейшего развития: добавления новых типов атак, звукового сопровождения, системы сохранения рекордов. Проект также может использоваться в учебных целях при изучении C++ и игрового программирования.

\end{document}